\chapter{The run model of the eraser}
\label{chap:run}

MetaCoq's certified eraser is a Rocq function over an inductive syntax, so its correctness proof
is an induction over that function's own equations [S, Sec.~7.3]. The Lean eraser is nothing of
the sort. It is production elaborator code: eighteen mutually recursive functions in a five-layer
monad stack, totalised by \code{partial\_fixpoint} rather than by a structural measure, calling
into \code{Lean.Meta}, the LCNF compiler front end and a kernel-backed relevance oracle, and
panicking --- successfully --- on shapes it does not expect. Before one can even \emph{state} that
its output is related to its input by \Erases, one needs a mathematical model of what ``the eraser
ran'' means. This chapter is that model: the file \code{LeanToLambdaBox/ErasureRun.lean}, a
deliberately \emph{model-free} library in which nothing mentions \Erases, lean4lean's \TrExprS{} or
\code{VEnv}, or the \lbox{} semantics. It is a program logic for one monad; everything that gives
the eraser a meaning is built on top of it in Chapters~\ref{chap:refinement}
and~\ref{chap:coldstart}.

The central design decision is stated in the file's own module docstring. Since \code{EraseM} is
\code{StateT ErasureState (ReaderT ErasureContext CoreM)} and \code{CoreM} unfolds to
\code{ReaderT Core.Context (StateRefT' IO.RealWorld Core.State (EIO Exception))}, an
\code{x : EraseM alpha} is \emph{run} by applying it to five arguments, and ``the run of \code{x}
succeeded'' is the bare equation \code{x s ctx cctx ref w = .ok (r, s') w'}. That spine is
deliberately not wrapped in a \code{def}, so that \code{rw}, \code{cases} and keyed matching see
the same head symbols the elaborated \code{do}-blocks produce. Four layers sit on top of it. The
\emph{run algebra} turns a successful \code{do}-block into a chain of successful steps and proves
--- rather than assumes --- that lifted actions leave the erasure state alone. The
\emph{admissibility toolkit} shows the run-ok motive to be \code{Lean.Order.admissible}, which is
what \code{partial\_fixpoint}'s generated induction principle demands of each of its eighteen
motives. The \emph{approximation toolkit} repairs a gap that fixpoint induction leaves open: the
principle hands each step an arbitrary point of the CCPO, so a conclusion about \emph{the} block
the shipping eraser builds needs an extra conjunct ``this function is below the shipping one''
carried through. And a small \emph{Hoare logic} covers exactly the loop forms the elaborated eraser
produces, including the parallel \code{for x in xs, y in ys} shape whose accumulator silently
threads a stream state.

The second half is the \emph{registration path}. Where MetaCoq relates a whole Rocq environment to
a \lbox{} environment by a separate judgement before any term is erased, the Lean eraser has no
environment pass at all: it builds the \lbox{} declarations \emph{during} the term walk,
demand-driven and memoised in the erasure state. That construction is modelled by \emph{proved}
state deltas. No node in this chapter carries a trust edge: the measured axiom footprint of every
declaration cited below is the standard one, and the file's only non-\code{LeanToLambdaBox} import
besides \code{Lean} is \code{Lean4Lean.Verify.NameGenerator}, a 28-line \code{sorry}-free module
used purely for the vocabulary \code{Lean.NameGenerator.Reserves} and its order.

\section{What a run is}
\label{sec:run:spelling}

\begin{definition}[A run of an \code{EraseM} computation]
\label{def:run-spelling}
\uses{def:Erasure-EraseM}
\code{EraseM} is \code{StateT ErasureState (ReaderT ErasureContext CoreM)}, and \code{CoreM} is
\code{ReaderT Core.Context (StateRefT' IO.RealWorld Core.State (EIO Exception))} with
\code{EIO eps = EST eps IO.RealWorld} and
\code{EST eps sigma alpha = Void sigma -> EST.Out eps sigma alpha}. An \code{x : EraseM alpha} is
therefore \emph{applied} to five arguments: the erasure state \code{s}, the erasure reader
\code{ctx}, the core reader \code{cctx : Core.Context}, the reference cell
\code{ref : ST.Ref IO.RealWorld Core.State} and the world token \code{w : Void IO.RealWorld}. We
say that \textbf{the run of \code{x} from \code{(s, ctx, cctx, ref, w)} succeeded with result
\code{r}, final
state \code{s'} and final world \code{w'}} when
\[
  \code{x s ctx cctx ref w} \;=\; \code{.ok (r, s') w'}.
\]
This raw equation, and not a wrapper definition, is the canonical hypothesis of every lemma below
and the canonical shape of every motive of every induction in
Chapters~\ref{chap:run}--\ref{chap:coldstart}: a wrapper constant would make lemma statements and
elaborated goals disagree about the head symbol.
\end{definition}

That this node has no Lean declaration of its own is precisely its content. ``The run of the
\emph{family}'' likewise means an application of
\code{Erasure.visitExpr.mutual\_fixpoint\_induct}, the principle \code{partial\_fixpoint} generates
for the eighteen-member \code{mutual} block; being auto-generated, it is named in prose only.

\begin{remark}
The fixpoint packs the eighteen members into one \code{PProd} chain
\code{Erasure.visitExpr.mutual}, whose slot order is \code{visitExpr}, \code{visitLiteral},
\code{visitConstructor}, \code{visitConst}, \code{get\_constant\_kername}, \code{visitMutual},
\code{visitAppArgs}, \code{visitLet}, \code{visitLambda}, \code{visitProj}, \code{visitApp},
\code{visitConstApp}, \code{visitCtorEta}, \code{visitCtorEtaGo}, \code{visitCasesEta},
\code{visitCasesEtaGo}, \code{visitCases}, \code{visitAlt}. This is \emph{not} the source order of
the \code{mutual} block in \code{Erasure.lean} --- \code{visitLambda} is third in the source and
ninth in the tuple, \code{visitConstructor} thirteenth in the source and third in the tuple --- and
every \code{motive\_1} \dots{} \code{motive\_18} numbering in the project follows the tuple order.
\end{remark}

\section{The run algebra}
\label{sec:run:algebra}

\begin{lemma}[Inversion for a successful bind]
\label{lem:Erasure-run_bind_ok}
\lean{Erasure.run_bind, Erasure.run_bind_ok, Erasure.run_bind_ne_ok}
\leanok
\uses{def:run-spelling}
Running a bind runs its first action and, on success, feeds the value \emph{together with} the new
state and the new world token to the continuation; on failure it propagates the error.
Consequently a bind's run succeeds exactly when there is a successful intermediate run:
\code{(x >>= f) s ctx cctx ref w = .ok (b, s') w'} if and only if there are \code{a}, \code{s\_1},
\code{w\_1} with \code{x s ctx cctx ref w = .ok (a, s\_1) w\_1} and
\code{f a s\_1 ctx cctx ref w\_1 = .ok (b, s') w'}. The three Lean declarations are the equation,
the inversion (an \code{iff}), and a corollary of its contrapositive (assuming \code{f} fails from
every possible intermediate result, not just the one \code{x} actually produces), which discharges
every error path of the eraser.
\end{lemma}

\begin{proof}
\leanok
\uses{def:run-spelling}
\code{run\_bind} is the one place where \code{do}-notation's match compilation is fought by hand:
the goal is restated at the \code{EST} layer, \code{EST.bind} unfolded, and a case split on the
first action's outcome closes both branches. Rewriting with it and casing again gives the
inversion, whose right-to-left direction is what makes it usable against opaque actions, where one
learns only that some intermediate result exists.
\end{proof}

Nearly every proof in this chapter and in Chapters~\ref{chap:refinement}--\ref{chap:coldstart}
opens by rewriting a run hypothesis with this inversion.

\begin{lemma}[The primitive runs]
\label{lem:Erasure-run_pure}
\lean{Erasure.run_pure, Erasure.run_get, Erasure.run_set, Erasure.run_modify,
Erasure.run_modifyGet, Erasure.run_read, Erasure.run_read_bind, Erasure.run_withReader,
Erasure.run_throw, Erasure.run_throw_ne_ok, Erasure.run_throwError_ne_ok}
\leanok
\uses{def:run-spelling}
\code{pure a} succeeds immediately with \code{a}, state and world unchanged; \code{get} returns the
state; \code{set s\_0} installs it; \code{modify g} replaces it by \code{g s}; \code{modifyGet g}
does both; \code{read} returns the erasure reader; \code{withReader f x} runs \code{x} under
\code{f ctx}; and \code{throw e} and \code{throwError msg} never return \code{.ok}. The chained
form \code{run\_read\_bind}, reducing \code{(read >>= f) s ctx cctx ref w} to
\code{f ctx s ctx cctx ref w}, steps the very first line of \code{visitExpr}, which reads the
reader's level parameters before invoking the relevance oracle.
\end{lemma}

\begin{proof}
\leanok
\uses{lem:Erasure-run_bind_ok}
All but three are \code{rfl}: the stack's \code{pure}, \code{get}, \code{set}, \code{modify},
\code{modifyGet}, \code{read}, \code{withReader} and \code{throw} reduce definitionally once the run
is applied to all five arguments. \code{run\_read\_bind} composes the bind equation with
\code{run\_read};
\code{run\_throw\_ne\_ok} is a \code{nomatch}; and \code{run\_throwError\_ne\_ok} unfolds
\code{Lean.throwError} into its \code{getRef} and \code{addMessageContext} binds, treats those as
opaque, and ends in \code{run\_throw\_ne\_ok}.
\end{proof}

\begin{remark}
The docstring of \code{run\_read} records that the inner \code{Core.Context} layer is \emph{not}
readable at \code{EraseM} --- there is no \code{MonadReaderOf Core.Context EraseM} instance --- and
is reached only through lifted \code{CoreM} actions.
\end{remark}

\begin{lemma}[Lifted actions are state-transparent]
\label{lem:Erasure-run_liftCoreM_ok}
\lean{Erasure.run_liftCoreM, Erasure.run_liftCoreM_ok, Erasure.run_liftCoreM_state,
Erasure.run_liftMetaM, Erasure.run_liftMetaM_ok, Erasure.run_liftMetaM_state}
\leanok
\uses{def:run-spelling}
An action lifted from \code{CoreM} cannot touch the erasure state: running
\code{(liftM x : EraseM alpha)} is running \code{x cctx ref w} with the erasure state carried
through, so a successful lifted run has final state equal to the entry state and moves only the
world token. The same holds for \code{Erasure.liftMetaM x}, which runs a \code{MetaM} action in
\code{CoreM} under the local context taken from the erasure reader. Each comes in three forms: the
raw equation, the \code{iff}-inversion and the state corollary.
\end{lemma}

\begin{proof}
\leanok
\uses{lem:Erasure-run_bind_ok, lem:Erasure-run_pure}
Both lifts unfold to a bind of a \code{CoreM} action with a \code{pure} that re-pairs the result
with the \emph{incoming} state; running that bind with the bind equation and \code{run\_pure} gives
the result, from which the inversion and the corollary are immediate.
\end{proof}

This makes \code{getEnv}, \code{mkFreshFVarId}, \code{logInfo}, \code{Meta.isInstance} and
\code{Compiler.LCNF.getDeclInfo?} state-transparent \emph{without any assumption}; the clauses the
world-indexed bundle of Chapter~\ref{chap:coldstart} still spends on them exist only because such a
primitive, while leaving the state alone, advances the world.

\begin{lemma}[A panic succeeds]
\label{lem:Erasure-run_panic}
\lean{Erasure.run_panic, Erasure.run_panicWithPosWithDecl}
\leanok
\uses{def:run-spelling}
At \code{EraseM} a \code{panic!} does \textbf{not} throw. With the \code{Inhabited} instance that
\code{panic!} and \code{unreachable!} pick up at type \code{EraseM alpha}, a panic returns
\code{default : alpha} and leaves state and world alone:
\code{(panic msg : EraseM alpha) s ctx cctx ref w = .ok (default, s) w}, and likewise for the
elaborated \code{panicWithPosWithDecl} form the macros actually produce. Every ``impossible''
branch of the eraser is therefore a reachable, \emph{successful} branch of the model. Both Lean
statements are \code{rfl}.
\end{lemma}

\begin{proof}
\leanok
Both are \code{rfl}: once the run is applied to all five arguments, the \code{Inhabited} instance
that the monad stack supplies reduces \code{panic} and its elaborated \code{panicWithPosWithDecl}
form to \code{pure default}.
\end{proof}

\begin{remark}
Three shipping consequences follow, and the verification models all three rather than excluding
them by hypothesis: \code{addAxiom}'s duplicate guard is a no-op that still registers a second
declaration entry (Theorem~\ref{thm:Erasure-run_addAxiom_ok}); every \emph{destructuring}
binder-helper lemma carries a
``the result was \code{default}'' disjunct, \ie{} the eraser can silently emit a default \lbox{}
term on an unexpected shape (Lemma~\ref{lem:Erasure-run_withLocalDecl_ok}); and
\code{register\_inductive}'s \code{unreachable!} arms fall through successfully. These are findings
about the shipping code, collected in Chapter~\ref{chap:trust}.
\end{remark}

\section{Admissibility: the run-ok motive}
\label{sec:run:adm}

\begin{lemma}[Admissibility at the \code{EST} leaf]
\label{lem:Erasure-est_admissible_ok}
\lean{Erasure.est_admissible_ok, Erasure.est_admissible_ok_pair}
\leanok
\uses{def:run-spelling}
At the bottom layer of the transformer stack the property ``on a successful run, \code{Q} holds''
is \code{Lean.Order.admissible}: for any \code{Q : Void sigma -> alpha -> Void sigma -> Prop}, the
predicate \code{fun x => forall w a w', x w = .ok a w' -> Q w a w'} on \code{EST eps sigma alpha}
is closed under suprema of chains. A second variant, with the result split as a pair, is the shape
produced after running the state-transformer layer.
\end{lemma}

\begin{proof}
\leanok
The proof works in the pointwise CCPO of flat orders over \code{EST.bot}, to which \code{EST}'s own
instance is definitionally equal; the instance must be supplied explicitly, since it is not
synthesised under the world binder. \code{admissible\_pi\_apply} and \code{admissible\_pi} peel the
quantifiers, and \code{admissible\_flatOrder} reduces the goal to the motive at bottom, which holds
vacuously because \code{EST.bot} is an \code{.error}.
\end{proof}

\begin{lemma}[The canonical run-ok motive is admissible]
\label{lem:Erasure-eraseM_admissible_ok}
\lean{Erasure.eraseM_admissible_ok}
\leanok
\uses{def:run-spelling, def:Erasure-EraseM}
For any predicate \code{Q} of a run's inputs and outputs, the property ``whenever this
computation's run succeeds, \code{Q} holds of the inputs and the outputs'' is an admissible
property of \code{EraseM tau}, and hence a legal motive for the fixpoint induction that
\code{partial\_fixpoint} generates. The Lean file adds five curried variants, of arity one to five,
for a member that takes that many arguments before returning an \code{EraseM} computation; between
them they cover all eighteen members. Their names carry Unicode subscripts and are therefore
described here rather than cited.
\end{lemma}

\begin{proof}
\leanok
\uses{lem:Erasure-est_admissible_ok}
Four applications of \code{admissible\_pi\_apply}, one per transformer layer, each with its
per-layer predicate supplied explicitly --- a bare \code{apply} fails higher-order unification ---
reduce the goal to the paired leaf lemma. Each arity variant adds one further application and
inherits from the previous one; arities four and five need a raised instance-synthesis budget.
\end{proof}

This single theorem is why the whole verification is phrased in run-ok form; the docstrings record
which members each arity serves.

\section{Approximation: staying below the fixpoint}
\label{sec:run:approx}

\begin{lemma}[Unsealing the eighteen slots]
\label{lem:Erasure-visitExpr_eq_mutual}
\lean{Erasure.visitExpr_eq_mutual, Erasure.visitLiteral_eq_mutual,
Erasure.visitConstructor_eq_mutual, Erasure.visitConst_eq_mutual,
Erasure.get_constant_kername_eq_mutual, Erasure.visitMutual_eq_mutual,
Erasure.visitAppArgs_eq_mutual, Erasure.visitLet_eq_mutual, Erasure.visitLambda_eq_mutual,
Erasure.visitProj_eq_mutual, Erasure.visitApp_eq_mutual, Erasure.visitConstApp_eq_mutual,
Erasure.visitCtorEta_eq_mutual, Erasure.visitCtorEtaGo_eq_mutual,
Erasure.visitCasesEta_eq_mutual, Erasure.visitCasesEtaGo_eq_mutual,
Erasure.visitCases_eq_mutual, Erasure.visitAlt_eq_mutual}
\leanok
\uses{def:run-spelling, def:Erasure-visitExpr}
\code{partial\_fixpoint} packs the mutual block into a single nested product
\code{Erasure.visitExpr.mutual} and seals each projection behind an irreducible definition. The
eighteen equations identify each named member with its slot --- \code{visitExpr} with the first
component, down to \code{visitAlt} with the last --- so that named constants and tuple slots are
interchangeable in lemma statements. Each Lean statement is \code{rfl} under an \code{unseal} of
the member it concerns.
\end{lemma}

\begin{proof}
\leanok
Each equation is \code{rfl}: unsealing the named member exposes the projection out of
\code{Erasure.visitExpr.mutual} that \code{partial\_fixpoint} defined it to be.
\end{proof}

\begin{lemma}[Below the fixpoint, a successful run is the fixpoint's run]
\label{lem:Erasure-run_ok_of_le}
\lean{Erasure.approx_rfl, Erasure.run_ok_of_le}
\leanok
\uses{def:run-spelling}
Let \code{x} and \code{y} in \code{EraseM tau} with \code{x} below \code{y} in the order
\code{partial\_fixpoint} uses. Then every successful run of \code{x} is verbatim a run of \code{y}:
same result, same final state, same final world token. The converse fails, which is why the
conjunct carried through the fixpoint induction must be the order itself and not an agreement of
successful runs. A one-argument variant, whose name carries a Unicode subscript, states the same
for a family member applied to an argument.
\end{lemma}

\begin{proof}
\leanok
The order on \code{EraseM tau} is pointwise down to the flat order over \code{EST.bot}, whose only
nontrivial relation is ``bottom is below anything''. Applying the hypothesis at the five run
arguments and casing on the resulting relation leaves two cases: in the bottom case the assumed
\code{.ok} equation is absurd, since \code{EST.bot} is an \code{.error}; in the reflexivity case
the two runs are literally the same term.
\end{proof}

This is the transport that lets a proof performed at the \emph{abstract} eraser handed out by
fixpoint induction be read at the \emph{shipping} \code{Erasure.visitExpr}.

\begin{lemma}[The extra conjunct costs nothing]
\label{lem:Erasure-fix_step_le}
\lean{Erasure.admissible_and_le, Erasure.fix_step_le}
\leanok
\uses{lem:Erasure-run_ok_of_le}
Adding the conjunct ``\dots{} and this function is below \code{c}'' to an admissible predicate
leaves it admissible; and one step of a monotone functional applied below its own fixpoint lands
below the fixpoint again. Together these discharge, for free, both the admissibility obligation and
the step obligation that the extra conjunct creates at each of the eighteen motives.
\end{lemma}

\begin{proof}
\leanok
\uses{lem:Erasure-run_ok_of_le}
The first is an application of \code{CCPO.csup\_le}: a chain that stays below \code{c} has its
supremum below \code{c}. The second is \code{Lean.Order.fix\_eq} read as an inequality ---
monotonicity moves the hypothesis under the functional, and \code{fix\_eq} turns one unfolding back
into the fixpoint.
\end{proof}

\begin{lemma}[The eighteen order conjuncts, packed]
\label{lem:Erasure-mutual_le_of}
\lean{Erasure.mutual_le_of}
\leanok
\uses{lem:Erasure-visitExpr_eq_mutual, def:Erasure-visitExpr}
Given eighteen functions of the family's eighteen signatures, each below the corresponding shipping
member, the tuple they form is below \code{Erasure.visitExpr.mutual}. Because the product order is
componentwise this is an eighteen-fold conjunction modulo the unsealing equations; composed with
Lemma~\ref{lem:Erasure-fix_step_le} at the family's own generated monotonicity proof, it discharges
the step obligation of the order conjunct at every one of the eighteen motives with a single
projection.
\end{lemma}

\begin{proof}
\leanok
\uses{lem:Erasure-visitExpr_eq_mutual}
Eighteen nested anonymous constructors, each rewriting the named member to its slot by the
corresponding unsealing equation. The only difficulty is the instance-synthesis budget:
\code{PartialOrder (EraseM tau)} is five transformer layers deep, so the default search size is
exhausted long before the eighteenth slot and is raised explicitly.
\end{proof}

\section{Hoare rules for the loops the eraser uses}
\label{sec:run:loops}

\begin{lemma}[Invariant rules for the plain loops]
\label{lem:Erasure-run_list_forIn_ok}
\lean{Erasure.run_list_forIn'_ok, Erasure.run_list_forIn_ok, Erasure.run_array_forIn_ok}
\leanok
\uses{def:run-spelling}
For an invariant \code{P} of the accumulator, the state and the world: if \code{P} holds initially
and is preserved by every successful body run --- whether the body yields or finishes early ---
then it holds of the result of a successful run of the loop. Three variants cover \code{forIn'}
over a list (with the membership proof available to the body), \code{forIn} over a list, and
\code{forIn} over an array. The array rule also covers the \emph{parallel}
\code{for x in xs, y in ys} shape, whose accumulator silently threads a stream state.
\end{lemma}

\begin{proof}
\leanok
\uses{lem:Erasure-run_bind_ok, lem:Erasure-run_pure}
Induction on the list, stepping the loop body with the bind inversion; the array rule rewrites
through \code{Array.forIn\_toList} to the list rule.
\end{proof}

\begin{lemma}[Prefix-indexed loop rules]
\label{lem:Erasure-run_list_forIn_ok-prime}
\lean{Erasure.run_list_forIn_ok', Erasure.run_array_forIn_ok', Erasure.run_list_foldlM_ok,
Erasure.run_array_foldlM_ok, Erasure.run_list_forIn_ok'_go, Erasure.run_list_foldlM_ok_go}
\leanok
\uses{lem:Erasure-run_list_forIn_ok}
In the prefix-indexed rules the invariant is indexed by the \emph{processed prefix}, so the step
hypothesis is handed the decomposition \code{L = pre ++ x :: post} and therefore knows which
element it is at and at which position; the early-exit arm must derive a contradiction rather than
be accommodated, and the conclusion is the invariant at the whole list. Four instances are provided
--- \code{forIn} over a list and over an array, \code{foldlM} over a list and over an array ---
each derived from an auxiliary induction the Lean file states separately.
\end{lemma}

\begin{proof}
\leanok
\uses{lem:Erasure-run_bind_ok}
Each auxiliary induction generalises over the pair of remaining list and processed prefix, so that
the list recursion goes through; the array forms rewrite through \code{Array.forIn\_toList} and
\code{Array.foldlM\_toList}.
\end{proof}

This is what a loop filling one output slot per input needs: \code{visitCases}' parallel
alternatives loop elaborates its two stream-exhaustion arms into early exits, which must be
\emph{refuted}, and the position index ties an output alternative to its input alternative
(Chapter~\ref{chap:refinement}).

\begin{lemma}[The prefix-indexed rule for \code{mapM}]
\label{lem:Erasure-run_list_mapM_ok}
\lean{Erasure.run_list_mapM_ok, Erasure.run_list_mapM_ok_go}
\leanok
\uses{lem:Erasure-run_list_forIn_ok-prime}
For an invariant seeing both the processed prefix and the outputs produced so far, and threading
state and world together: if it holds of the empty prefix and the empty output, and every
successful body run extends it by one element and one output, then it holds of the whole list and
the whole output list after a successful run of \code{List.mapM f L}.
\end{lemma}

\begin{proof}
\leanok
\uses{lem:Erasure-run_bind_ok, lem:Erasure-run_pure}
The auxiliary induction generalises over the remaining list and the processed prefix and
additionally accounts for the reversed accumulator of \code{List.mapM.loop}; the rule follows by
instantiating the prefix at the empty list.
\end{proof}

This is the most-used loop rule downstream --- \code{visitMutual}'s sibling loop, the two loops of
\code{register\_inductive}, the fresh-identifier loop and \code{mkAlt}'s name loop all go through
it --- and it threads \emph{state and world together}, which is what later makes distinctness of
freshly minted identifiers provable (Theorem~\ref{thm:Erasure-run_mkFreshFVarId_list}). A section of
worked examples documents how the rules compose, two of them on the parallel \code{for} shape.

\section{A scale check on the whole family}
\label{sec:run:scale}

\begin{theorem}[A scale check on the whole family]
\label{thm:Erasure-visitExpr_run_shape}
\lean{Erasure.visitExpr_run_shape}
\leanok
\uses{def:run-spelling, def:Erasure-visitExpr, def:LBTerm}
An eighteen-fold conjunction, one clause per family member, each in the canonical run-ok form, of
which four carry real content: a successful \code{visitExpr} run on a lambda input returns $\Box$
or a \lbox{} lambda; a successful \code{visitConst} run on a constant input returns a free variable
or a constant; a successful \code{visitAppArgs} run with a non-empty argument array returns an
application; and a successful \code{visitLambda} run on a lambda input returns a lambda. The
remaining fourteen clauses conclude \code{True}.
\end{theorem}

\begin{proof}
\leanok
\uses{lem:Erasure-eraseM_admissible_ok, lem:Erasure-run_bind_ok, lem:Erasure-run_pure,
lem:Erasure-run_list_forIn_ok-prime}
By the family's generated fixpoint induction with all eighteen motives in run-ok form. The eighteen
admissibility obligations discharge from the arity variants of
Lemma~\ref{lem:Erasure-eraseM_admissible_ok}; each step goal unfolds the member, steps its binds
with the run algebra, and either appeals to a sibling induction hypothesis or concludes
\code{True}. The \code{visitAppArgs} clause goes through the array \code{foldlM} rule.
\end{proof}

\begin{remark}
The theorem's own docstring calls it a \emph{scale check}: it confirms that the eighteen
admissibility obligations discharge, that the step goals are tractable with the run lemmas, and
that elaboration time is acceptable. It has no downstream consumer. The load-bearing
eighteen-motive inductions are \code{visitExpr\_shapeW} (Chapter~\ref{chap:coldstart}) and the
refinement induction of Chapter~\ref{chap:refinement}; this one is not a step of the capstone.
\end{remark}

\section{State deltas: what a run may do to the registry}
\label{sec:run:state}

\begin{lemma}[The ambient primitives do not touch the erasure state]
\label{lem:Erasure-run_getConstInfo_state}
\lean{Erasure.run_monadRefWithRef, Erasure.run_logInfo_state, Erasure.run_getEnv_state,
Erasure.run_mkFreshFVarId_state, Erasure.run_getConstInfo_state}
\leanok
\uses{def:run-spelling}
A successful run of \code{logInfo}, \code{getEnv}, \code{mkFreshFVarId} or \code{getConstInfo}
leaves the erasure state unchanged. The first three are corollaries of lifting
(Lemma~\ref{lem:Erasure-run_liftCoreM_ok}); the fourth is not --- \code{getConstInfo} elaborates
\emph{at} \code{EraseM} rather than as a literal lift --- and is proved outright.
\end{lemma}

\begin{proof}
\leanok
\uses{lem:Erasure-run_bind_ok, lem:Erasure-run_liftCoreM_ok, lem:Erasure-run_pure}
\code{Lean.getConstInfo} unfolds to \code{getEnv} followed by a match on
\code{Environment.find?}; the success branch is a \code{pure}, and the failure branch is refuted by
chasing \code{throwUnknownConstant} through \code{throwUnknownConstantAt},
\code{throwUnknownIdentifierAt}, \code{throwErrorAt} and \code{withRef} down to the fact that a
thrown error never returns \code{.ok}. \code{run\_monadRefWithRef} is the small \code{rfl} that
pushes through \code{withRef}.
\end{proof}

The section docstring singles this lemma out as load-bearing: because it is a \emph{theorem}, no
hypothesis bundle in the development has to assume that an environment lookup is state-transparent.

\begin{lemma}[An \code{EraseM} lookup is a \code{CoreM} lookup]
\label{lem:Erasure-pass_getConstInfo_core}
\lean{Erasure.pass_core_bind, Erasure.pass_getConstInfo_core}
\leanok
\uses{def:run-spelling}
A successful run of \code{Lean.getConstInfo} at \code{EraseM} yields a successful run of
\code{Lean.getConstInfo} at \code{CoreM}, on the same name, with the same result and the same world
tokens. The two elaborations are not definitionally equal, so this bridge is proved;
\code{pass\_core\_bind} is the bind equation one layer down, at \code{CoreM}.
\end{lemma}

\begin{proof}
\leanok
\uses{lem:Erasure-run_bind_ok, lem:Erasure-run_getConstInfo_state, lem:Erasure-run_liftCoreM_ok,
lem:Erasure-run_pure}
Both sides unfold to \code{getEnv} plus the same \code{Environment.find?} match; the error branches
are refuted exactly as in Lemma~\ref{lem:Erasure-run_getConstInfo_state}, and the success branches
are matched by \code{pass\_core\_bind}.
\end{proof}

This lets an adequacy specification stated at \code{CoreM} --- the lookup-adequacy clause of the
specification bundle of Chapter~\ref{chap:eraser} --- be applied to the eraser's own calls. It is a
bridge between two Lean monads, not a bridge to lean4lean.

\begin{definition}[Canonical constant registry]
\label{def:Erasure-CanonicalConstants}
\lean{Erasure.CanonicalConstants}
\leanok
\uses{def:Erasure-EraseM, def:Kername}
The erasure state's constant registry is \textbf{canonical} when every constant it knows is
registered under its own canonical kername: whenever the registry maps \code{n} to \code{k},
\code{k} is \code{toKername n}. This is the one value-level invariant carried across a growing
state.
\end{definition}

\begin{definition}[The modelled registration states]
\label{def:Erasure-addAxiomState}
\lean{Erasure.addAxiomState, Erasure.mutualBlockKn, Erasure.registerIndState}
\leanok
\uses{def:Erasure-CanonicalConstants, def:GlobalDeclarations, def:Erasure-EraseM,
def:OneInductiveBody}
\code{addAxiomState n s} is the exact post-state of registering \code{n} as an axiom: the registry
gains \code{n} mapped to \code{toKername n}, and the global declaration list gains one value-less
constant declaration under that kername. \code{mutualBlockKn ii} is the kername the eraser mints
for a mutual inductive \emph{block}, the root kername of the concatenation of the block members'
names. \code{registerIndState ii bodies s} is the post-state of the block registration, which
prepends one inductive declaration carrying the block's parameter count and its bodies.
\end{definition}

\begin{definition}[The two constant-registration states]
\label{def:Erasure-nonrecConstState}
\lean{Erasure.nonrecConstState, Erasure.recConstState, Erasure.recConstStep,
Erasure.recConstState_eq}
\leanok
\uses{def:Erasure-CanonicalConstants, def:LBTerm, def:FixDef}
\code{nonrecConstState n t s} is the post-state of the non-recursive exit: the registry gains
\code{n} under its canonical kername and the declaration list gains a constant declaration whose
body is \code{t}. \code{recConstState names defs s} is the post-state of the block exit: a fold
over the block's names paired with their indices, registering each sibling under its own canonical
kername with body the fixpoint block \code{defs} selected at that index. Its single step,
\code{recConstStep}, is exactly \code{nonrecConstState} at a fixpoint body, named separately so
that the fold induction has something to generalise over; \code{recConstState\_eq} is the
\code{rfl} relating the two.
\end{definition}

\begin{definition}[Axiom-only state extension]
\label{def:Erasure-ConstExt}
\lean{Erasure.ConstExt, Erasure.AxiomExt, Erasure.ConstExt.rfl', Erasure.ConstExt.of_same,
Erasure.ConstExt.trans, Erasure.AxiomExt.rfl', Erasure.AxiomExt.trans, Erasure.AxiomExt.addAxiom}
\leanok
\uses{def:Erasure-CanonicalConstants, def:Erasure-addAxiomState, def:GlobalDeclarations}
\code{ConstExt s s'} says that \code{s'} is an \emph{axiom-only} extension of \code{s}, in three
fields: canonicity of the constant registry is preserved; the registry's domain only grows; and the
global declaration list grows by a prefix consisting solely of value-less constant declarations,
each keyed by the canonical kername of a constant the extended registry knows. \code{AxiomExt}
extends \code{ConstExt} by one further field requiring that the inductive registry is untouched.
Both are reflexive and transitive, and the modelled post-state of an axiom registration really is
an \code{AxiomExt} extension.
\end{definition}

\begin{remark}
Recording the \emph{keys} of the prefix, and not merely the shape of its entries, is what lets a
downstream key-coverage invariant cross a \code{register\_inductive} call, whose cold branch emits
one axiom per externally implemented constructor. \code{AxiomExt} itself has no consumer outside
this file.
\end{remark}

\begin{definition}[State extension]
\label{def:Erasure-StateLe}
\lean{Erasure.StateLe, Erasure.StateLe.rfl', Erasure.StateLe.trans}
\leanok
\uses{def:Erasure-EraseM, def:GlobalDeclarations}
\code{StateLe s s'} says that the two registries only grew --- every name the constant registry
knew it still knows, and likewise for the inductive registry --- and that the global declaration
list only got prepended to. It is reflexive and transitive, and is the monotonicity a \emph{cold}
run has in place of a warm bridge's ``the state did not change''.
\end{definition}

\begin{definition}[The widened run conclusion]
\label{def:Erasure-RunConcl}
\lean{Erasure.RunConcl, Erasure.RunConcl.rfl', Erasure.RunConcl.of_eq, Erasure.RunConcl.trans}
\leanok
\uses{def:Erasure-StateLe, def:Erasure-CanonicalConstants}
\code{RunConcl s s'} bundles state extension with survival of canonicity, in two fields: the state
only grew, and if the constant registry was canonical before, it is canonical after. It is
reflexive, transitive, and holds whenever the state did not change at all.
\end{definition}

\code{RunConcl} is central: it is the shape every motive of the refinement induction of
Chapter~\ref{chap:refinement} concludes. Its second field is \emph{provable}, never assumed, because
every writer on the registration path inserts a constant under its own canonical kername ---
Lemma~\ref{lem:Erasure-runConcl_nonrecConstState} and
Theorem~\ref{thm:Erasure-run_register_inductive_runConcl} are exactly that.

\begin{definition}[A registry entry justified by the emitted block]
\label{def:Erasure-RegisteredBodyAt}
\lean{Erasure.RegisteredBodyAt, Erasure.RegisteredBodyAt.mono}
\leanok
\uses{def:Erasure-addAxiomState, def:OneInductiveBody, def:InductiveId}
The inductive-registry entry \code{rc} for a name \code{n} is \emph{justified by} the block
\code{bodies} emitted for \code{ii} when it names that block, points at a body of \code{bodies}
whose name is \code{n}, and its constructor-argument masks keep exactly as many arguments as that
body's constructors declare. Justification is stable under extending \code{bodies} on the right.
The predicate has no consumer outside this file; it is the content of the cold branch's registry
effect as stated by Theorem~\ref{thm:Erasure-run_register_inductive_cold_ok}.
\end{definition}

\begin{lemma}[A split of an indexed list names its position]
\label{lem:Erasure-zipIdx_split_snd}
\lean{Erasure.zipIdx_split_snd, Erasure.zipIdx_split_fst}
\leanok
If pairing a list with its indices splits as \code{pre ++ x :: post}, then the index component of
\code{x} is the length of \code{pre}, and the original list has \code{x}'s value at that position.
This is what turns the prefix decomposition handed out by the prefix-indexed loop rules into an
index into the original list.
\end{lemma}

\begin{proof}
\leanok
Read the element at position \code{pre.length} off both sides of the split equation and unfold the
index-pairing lemma for that position.
\end{proof}

\section{The registration path}
\label{sec:run:reg}

\begin{theorem}[The whole state effect of an axiom registration]
\label{thm:Erasure-run_addAxiom_ok}
\lean{Erasure.run_addAxiom_ok}
\leanok
\uses{def:Erasure-addAxiomState, def:Erasure-register_inductive}
Every successful run of \code{Erasure.addAxiom n} ends in the state \code{addAxiomState n s} and
leaves the world token unchanged --- \textbf{unconditionally}, on both branches of the function's
``already defined'' guard.
\end{theorem}

\begin{proof}
\leanok
\uses{lem:Erasure-run_bind_ok, lem:Erasure-run_pure, lem:Erasure-run_panic}
Step the \code{get}, then split on the guard. In the ``already present'' branch the shipping code
runs \code{panic!} and, having no \code{return}, falls straight through into the same
\code{modify}; because a panic \emph{succeeds} at \code{EraseM}
(Lemma~\ref{lem:Erasure-run_panic}), that branch reaches the modification too. Both branches
therefore end at \code{addAxiomState n s}, and neither moves the world.
\end{proof}

\begin{remark}
A finding about the shipping code, modelled rather than assumed away: the duplicate guard is a
no-op, so a second, duplicate entry is prepended to the declaration list.
\end{remark}

\begin{theorem}[The miss branch of \code{register\_inductive} is not state-preserving]
\label{thm:Erasure-run_register_inductive_cold_ok}
\lean{Erasure.run_register_inductive_cold_ok, Erasure.run_register_inductive_hit_ok}
\leanok
\uses{def:Erasure-ConstExt, def:Erasure-addAxiomState, def:Erasure-RegisteredBodyAt,
def:Erasure-register_inductive}
If the inductive registry already knows the block, a successful \code{register\_inductive} run
returns the stored entry and changes neither state nor world. If it does not, a successful run
produces a list of inductive bodies \code{bodies} and an intermediate state \code{sM} such that the
final state is \code{registerIndState ii bodies sM}, the result is the entry \code{sM} holds for
the block's own name, \code{bodies} has one entry per block member, \code{sM} is an axiom-only
extension of the entry state, the inductive registry only grew, and every inductive-registry entry
present in \code{sM} was either already present or is justified by the block just emitted.
\end{theorem}

\begin{proof}
\leanok
\uses{lem:Erasure-run_bind_ok, lem:Erasure-run_pure, lem:Erasure-run_list_mapM_ok,
lem:Erasure-run_getConstInfo_state, thm:Erasure-run_addAxiom_ok, lem:Erasure-run_panic,
def:Erasure-ConstExt, lem:Erasure-zipIdx_split_snd}
Step the registry test; the hit branch is a \code{get} followed by a \code{pure}. On the miss
branch the body maps over the block's members, so the prefix-indexed \code{mapM} rule carries an
invariant recording the output length, an axiom-only extension of the entry state, growth of the
inductive registry, and justification of every new entry. The inner constructor loop uses the same
rule again and contributes one axiom registration per externally implemented constructor,
discharged by Theorem~\ref{thm:Erasure-run_addAxiom_ok}; the position of each member comes from
Lemma~\ref{lem:Erasure-zipIdx_split_snd}. Finally one inductive declaration is prepended.
\end{proof}

\begin{remark}
This says plainly that any specification asserting that \code{register\_inductive} preserves the
state at an arbitrary entry state is \emph{false about the real function}, which is why every
registration premise in the development is guarded by a registry test. The file also contains a
companion that constructs a hit-branch run out of nothing --- a refutation device for unguarded
premises --- which has no consumer in the current tree, and whose docstring cites a refutation that
no longer exists.
\end{remark}

\begin{theorem}[What the cold branch leaves in the inductive registry]
\label{thm:Erasure-run_register_inductive_cold_entries}
\lean{Erasure.run_register_inductive_cold_entries}
\leanok
\uses{def:Erasure-addAxiomState, def:Erasure-RegisteredBodyAt, def:Erasure-ErasureConfig}
Fix an abstract lookup predicate \code{Ci} --- ``what a successful environment lookup reports'' ---
and a reading \code{gw} of the name generator out of the world token, and assume that every lookup
during the run reports \code{Ci}, and that every lookup, environment read and log during the run
advances \code{gw} monotonically.
Then, on the miss branch and with constructor-argument pruning switched off, the generator only
advances, and every inductive-registry entry present afterwards was either already present or names
the block identifier minted from the block's member list at the member's own position, with the
member itself reported by \code{Ci} as an inductive, and carries, for each constructor, an all-keep
mask of exactly that constructor's field count.
\end{theorem}

\begin{proof}
\leanok
\uses{thm:Erasure-run_register_inductive_cold_ok, lem:Erasure-run_list_mapM_ok,
lem:Erasure-run_getConstInfo_state, lem:Erasure-run_bind_ok}
The same walk as the previous theorem, with the loop invariant carrying, in addition, the abstract
lookup predicate at every member and constructor and the generator bound. The pruning flag is what
fixes the mask: with pruning off, the mask computed per constructor is a replication of ``keep'' at
the constructor's field count.
\end{proof}

The predicate \code{Ci} is what keeps this chapter model-free: the model reading instantiates it at
``the source environment holds this constant information''. It is also the cleanest place to observe
that the pipeline \emph{declines} the constructor-argument pruning that ConCert-style typed
extraction performs.

\begin{theorem}[The run conclusion of a registration]
\label{thm:Erasure-run_register_inductive_runConcl}
\lean{Erasure.run_register_inductive_runConcl, Erasure.run_register_inductive_entries,
Erasure.run_register_inductive_gdeclsConst}
\leanok
\uses{def:Erasure-RunConcl, def:Erasure-register_inductive}
Every successful \code{register\_inductive} run satisfies \code{RunConcl}: the state only grew and
canonicity survived. The entry characterisation of the previous theorem extends to both branches,
the hit branch contributing nothing. Moreover the call never records a constant \emph{body}: every
constant declaration carrying a body that is present afterwards was already present before.
\end{theorem}

\begin{proof}
\leanok
\uses{thm:Erasure-run_register_inductive_cold_ok, thm:Erasure-run_register_inductive_cold_entries,
def:Erasure-ConstExt, def:Erasure-StateLe, def:Erasure-RunConcl}
Case on the registry test. The hit branch preserves the state outright, so \code{RunConcl} follows
from its reflexivity. The miss branch supplies an axiom-only extension whose declaration prefix,
extended by the single inductive declaration the branch prepends, is the prefix \code{StateLe}
needs, and whose canonicity clause is the one \code{RunConcl} needs. The body-freedom statement
inspects that prefix: every entry it contains is either the block declaration or a value-less
constant declaration.
\end{proof}

Nothing is assumed here: this is the whole state effect of the call, proved. The body-freedom
corollary has no consumer in the current tree, but it is the fact that lets a ``the stored body
erases its source body'' record cross the call for free.

\begin{theorem}[The shape of an emitted fixpoint definition]
\label{thm:Erasure-run_mkDef_ok}
\lean{Erasure.run_mkDef_ok, Erasure.isLambda_foldl_toBvar, Erasure.run_mkDef_isLambda,
Erasure.run_mkDef_rarg}
\leanok
\uses{def:Erasure-mkLambda, def:toBvar, def:FixDef, def:atomValue}
A successful run of \code{Erasure.mkDef nm fixvarnames body} names the definition after the Lean
constant, closes the erased body over the block's fixpoint variables by a fold of de Bruijn
abstractions --- in reverse order, at levels given by position --- and touches neither state nor
world. Two consequences of that shape are stated separately: the emitted definition is
lambda-headed exactly when the erased body is, and its principal-argument index is the default
\code{0}, because \code{mkDef} never sets that field.
\end{theorem}

\begin{proof}
\leanok
\uses{lem:Erasure-run_bind_ok, lem:Erasure-run_list_forIn_ok-prime, lem:Erasure-run_pure}
The body is a single \code{for} loop over the reversed, index-paired fixpoint names, each iteration
reading the reader and abstracting one variable; the prefix-indexed \code{forIn} rule with the
invariant ``the accumulator is the fold over the processed prefix'' gives the equation, and the
early-exit arm is refuted because the body always yields. A de Bruijn abstraction never changes the
head constructor, which gives the head-shape corollary; the principal-argument corollary is
\code{rfl} after the loop is stepped.
\end{proof}

\begin{remark}
These are exactly the side conditions the target semantics' fixpoint rules need
(Chapter~\ref{chap:lambdabox}): a fixpoint definition must be lambda-headed, and its recursive
argument index must agree with what the unfolding rule reads. The principal-argument corollary has
no consumer in the current tree, and its docstring attributes it to a fixpoint rule of \Erases{}
that does not exist --- the erasure relation has no fixpoint constructor; the corresponding premise
lives on the lowering relation's fixpoint rules (Chapter~\ref{chap:lowering}).
\end{remark}

\begin{lemma}[A loop of state modifications is a fold]
\label{lem:Erasure-run_modify_forIn_ok}
\lean{Erasure.run_modify_forIn_ok}
\leanok
\uses{def:run-spelling}
A loop \code{for x in L do modify (g x)} ends at the left fold of \code{g} over \code{L} from the
entry state, with the world unchanged. This is the block-registration loop of \code{visitMutual}'s
recursive exit, and it is what makes \code{recConstState} the right model for that exit.
\end{lemma}

\begin{proof}
\leanok
\uses{lem:Erasure-run_list_forIn_ok-prime, lem:Erasure-run_pure}
The prefix-indexed \code{forIn} rule with the invariant ``the state is the fold over the processed
prefix''; the early-exit arm is refuted because the body always yields.
\end{proof}

\begin{lemma}[The memoising constant lookup]
\label{lem:Erasure-run_get_constant_kername_ok}
\lean{Erasure.run_get_constant_kername_ok}
\leanok
\uses{def:run-spelling, def:Erasure-visitMutual}
A successful run of \code{get\_constant\_kername n} is either a registry hit --- the registry
already maps \code{n} to the returned kername, and neither state nor world moved --- or a registry
miss followed by a successful \code{visitMutual n} run, after which the returned kername is what the
registry then holds for \code{n}. This is the demand-driven, memoised environment construction in
miniature; the lemma has no consumer in the current tree.
\end{lemma}

\begin{proof}
\leanok
\uses{lem:Erasure-run_bind_ok, lem:Erasure-run_pure}
Step the \code{get} and case on the registry lookup; the hit arm is a \code{pure}, the miss arm
binds \code{visitMutual} and re-reads the registry.
\end{proof}

\section{Hoare rules for \code{visitMutual}'s exits}
\label{sec:run:exits}

\begin{proposition}[\code{visitMutual} in Hoare form over its four exits]
\label{prop:Erasure-run_visitMutual_ok}
\lean{Erasure.run_visitMutual_ok}
\leanok
\uses{def:Erasure-addAxiomState, def:Erasure-nonrecConstState, def:Erasure-visitMutual,
def:Erasure-prepare_erasure}
Let \code{Q} be a predicate of the erasure state closed under: the inline bookkeeping entry; the
axiom-registration delta; a run of the preprocessing pass; a run of the shipping \code{visitExpr},
which must in addition yield a body satisfying two abstract output predicates; the non-recursive
constant registration, given those two predicates of the stored body; and the block registration,
given the block's length and the fact that each of its bodies is the abstraction closure of a good
erasure output. Then \code{Q} survives every successful \code{visitMutual} run.
\end{proposition}

\begin{proof}
\leanok
\uses{prop:Erasure-run_nonrec_exit_ok, lem:Erasure-run_bind_ok, lem:Erasure-run_liftCoreM_ok,
lem:Erasure-run_getConstInfo_state, thm:Erasure-run_addAxiom_ok, lem:Erasure-run_pure}
Step the compiler-declaration query and the environment read, both state-transparent, then split on
the inline-attribute test. The elaborated body's core is a match on three discriminants at once ---
whether the declaration has a value, whether it is externally implemented, and what the
configuration says about externals --- which the tactic \code{split} cannot handle and which is
therefore resolved by three nested case analyses. Each resulting arm is one of the four exits: the
axiom exit closes by Theorem~\ref{thm:Erasure-run_addAxiom_ok}, the two registering exits by the
rules of Proposition~\ref{prop:Erasure-run_nonrec_exit_ok}, and the inline prefix by its own rule.
\end{proof}

\begin{remark}
The clause about the preprocessing pass is a \emph{hypothesis} of this proposition, not a proved
fact: the pass's simplification branch runs a \code{Core} transformation \emph{at} \code{EraseM}, so
its state transparency does not follow from the lifting lemmas. The proposition is therefore
conditional on it. The clause \emph{is} discharged downstream, under the pinned configuration, by
\code{LeanToLambdaBox.run\_prepare\_erasure\_state} (Chapter~\ref{chap:coldstart}), which proves it
outright when the simplification flag is off; the case where the flag is on remains unhandled. The
Lean docstring still attributes the clause to a trust class that no longer exists in the tree.
\end{remark}

\begin{proposition}[The four exit rules, state form]
\label{prop:Erasure-run_nonrec_exit_ok}
\lean{Erasure.run_inline_tail_ok, Erasure.run_inline_prefix_ok, Erasure.run_nonrec_exit_ok,
Erasure.run_rec_exit_ok}
\leanok
\uses{def:Erasure-nonrecConstState, def:Erasure-addAxiomState, def:Erasure-mkLambda}
Four Hoare rules, one per shape \code{visitMutual} exits through, each with every boolean test, log
message and reader update left abstract. The inline \emph{tail} rule covers the post-registration
bookkeeping, which at most prepends one entry to the inline list; the inline \emph{prefix} rule
covers the attribute test that may precede the core and then runs the same continuation on both
branches. The non-recursive exit erases the preprocessed body under the declaration's reader update,
prepends the constant, and runs the inline tail. The recursive exit mints one fresh variable per
sibling, erases each sibling body through \code{mkDef} under the fixpoint-variable binding, and
prepends one registry entry per name; its closure hypothesis is handed the \emph{shape} of the block
being stored --- one definition per name, each definition's body the abstraction closure of a good
erasure output over the block's names.
\end{proposition}

\begin{proof}
\leanok
\uses{lem:Erasure-run_bind_ok, lem:Erasure-run_pure, lem:Erasure-run_list_mapM_ok,
thm:Erasure-run_mkDef_ok, lem:Erasure-run_modify_forIn_ok, lem:Erasure-run_getConstInfo_state}
Each rule steps its own \code{do}-block with the bind inversion, using the \code{mapM} rule for the
sibling loop, the definition-shape theorem for each sibling's \code{mkDef} call, and the fold rule
for the registration loop; the inline tail's two branches differ only in which message is logged.
\end{proof}

The eraser is kept \emph{abstract} in these rules because inside the family's fixpoint induction the
step goal for \code{visitMutual} mentions the fixpoint's abstract argument, not the shipping
\code{visitExpr}. And the recursive rule cannot simply demand closedness of an arbitrary block --- a
fixpoint block whose body mentions an out-of-range de Bruijn index is not closed --- so the caller
must be able to \emph{compute} that closedness from the per-definition shape.

\begin{proposition}[The four exit rules, world-indexed]
\label{prop:Erasure-run_nonrec_exit_ok-prime}
\lean{Erasure.run_inline_tail_ok', Erasure.run_inline_prefix_ok', Erasure.run_nonrec_exit_ok',
Erasure.run_rec_exit_ok'}
\leanok
\uses{prop:Erasure-run_nonrec_exit_ok, def:Erasure-nonrecConstState}
The same four rules for a predicate of the erasure state \textbf{and the world token}. Two changes
are forced by the extra index. First, every state-transparent primitive on the registration path ---
logging, the instance test, identifier minting, environment lookup --- leaves the state alone but
\emph{advances the world}, so each now needs its own preservation clause keyed on its own run.
Second, the erasure clause is keyed on the preprocessing run as well as on the erasure run, so that
a caller who must \emph{re-establish} an invariant at the erasure's entry state can see the run that
produced it.
\end{proposition}

\begin{proof}
\leanok
\uses{prop:Erasure-run_nonrec_exit_ok, lem:Erasure-run_bind_ok, thm:Erasure-run_mkDef_ok,
lem:Erasure-run_modify_forIn_ok, lem:Erasure-run_list_mapM_ok}
The same walks, with the invariant now carrying the world token; the two world-preserving steps ---
the definition builder and the registration loop --- need no clause, which is why the recursive rule
has fewer hypotheses than one might expect.
\end{proof}

The reason for the duplication is stated in the file: the bridge's motives conclude a
\emph{relation} between two states together with a \emph{bound on the name generator}, and the
generator is read out of the world token. A state-only predicate can express neither.

\begin{lemma}[Stepping past the inline prefix]
\label{lem:Erasure-run_inline_prefix_decomp-prime}
\lean{Erasure.run_inline_prefix_decomp'}
\leanok
\uses{def:run-spelling, def:Kername}
Instead of propagating an invariant through the inline-attribute prefix, this rule
\emph{decomposes} it: there are an intermediate state and world such that either nothing happened,
or one logging run advanced the world at unchanged state and the state gained one inline entry; and
the continuation ran from there.
\end{lemma}

\begin{proof}
\leanok
\uses{lem:Erasure-run_bind_ok, lem:Erasure-run_getConstInfo_state, lem:Erasure-run_liftCoreM_ok}
Split on the attribute test and step each branch; the only state write is the inline entry.
\end{proof}

This is the form a caller needs whose conclusion is \emph{false} at the call's entry and becomes true
only at the registration --- ``\code{n} is in the registry when this call returns'' --- so that no
single predicate is both establishable at entry and strong enough at exit. The environment step of
the refinement (Chapter~\ref{chap:refinement}) consumes it.

\begin{lemma}[Each registration delta is a run conclusion]
\label{lem:Erasure-runConcl_nonrecConstState}
\lean{Erasure.runConcl_inlinings, Erasure.runConcl_addAxiomState,
Erasure.runConcl_nonrecConstState, Erasure.runConcl_foldl_recConstStep,
Erasure.runConcl_recConstState, Erasure.nonrecConstState_get?, Erasure.addAxiomState_get?,
Erasure.foldl_recConstStep_get?, Erasure.recConstState_get?}
\leanok
\uses{def:Erasure-RunConcl, def:Erasure-nonrecConstState, def:Erasure-addAxiomState}
Each of the four state deltas the eraser performs is a \code{RunConcl} step: the inline entry writes
only the inline list; the axiom exit registers one constant under its own canonical kername and
prepends one value-less declaration; the non-recursive exit does the same with a body; and the block
registration, being a fold of the non-recursive step at fixpoint bodies, composes out of the
previous one. Four further lemmas record that each of these registrations really does put the name
it registers into the registry.
\end{lemma}

\begin{proof}
\leanok
\uses{def:Erasure-RunConcl, def:Erasure-nonrecConstState}
Each delta inserts at most one registry key, mapped to that key's canonical kername, and prepends at
most one declaration, so the three fields of state extension and the canonicity field are read off
directly; the block registration is an induction over the fold, using transitivity of
\code{RunConcl}.
\end{proof}

Together with Theorem~\ref{thm:Erasure-run_register_inductive_runConcl} these are the \emph{proof}
that the canonicity field of \code{RunConcl} is never assumed.

\begin{theorem}[The identifier-minting loop]
\label{thm:Erasure-run_mkFreshFVarId_list}
\lean{Erasure.run_mkFreshFVarId_list}
\leanok
\uses{def:run-spelling}
Assume a freshness specification for identifier minting: each successful run returns an identifier
not reserved by the generator before the run, reserved by the generator after it, with the generator
advanced monotonically, and reserved by a fixed ambient generator. Then a successful run of one
minting per name in a list returns exactly as many identifiers as there are names, pairwise
distinct, at an unchanged state, with the generator advanced, and with every returned identifier
reserved by the final generator, reserved by the ambient one, and outside any given list of
identifiers all of which the initial generator reserves.
\end{theorem}

\begin{proof}
\leanok
\uses{lem:Erasure-run_list_mapM_ok, lem:Erasure-run_getConstInfo_state}
One application of the \code{mapM} rule whose invariant chains state and world together.
Distinctness is the payoff of that chaining: at each step the new identifier is not reserved by the
generator read from the current world, while every earlier identifier is --- carried in the
invariant and transported along generator monotonicity --- so the new one differs from all of them.
Disjointness from the ambient list is the same argument run against the hypothesis that the initial
generator reserves that list.
\end{proof}

The generator vocabulary is lean4lean's, from a 28-line \code{sorry}-free module, so this inherits no
trust. It has no consumer in the current tree: part of the library, not a step of the capstone.

\begin{theorem}[The sibling loop, chained]
\label{thm:Erasure-run_rec_exit_siblings_chained}
\lean{Erasure.run_rec_exit_siblings_chained}
\leanok
\uses{def:run-spelling, def:Erasure-prepare_erasure, def:Erasure-mkLambda, def:FixDef}
For a per-sibling package \code{R} and a step hypothesis that is handed the four primitive runs
\emph{in run order and at the states and worlds they actually consume} --- the environment lookup,
the preprocessing pass under the sibling's own reader, the abstract eraser, and the definition
builder --- a successful run of \code{visitMutual}'s sibling loop produces exactly as many
definitions as there are names, carries the invariant from the loop's entry to its exit, and
returns, for every index, a definition at that index satisfying \code{R} for the name at that index.
\end{theorem}

\begin{proof}
\leanok
\uses{lem:Erasure-run_list_mapM_ok, lem:Erasure-run_bind_ok, lem:Erasure-run_pure}
One application of the \code{mapM} rule whose invariant carries the output length, the caller's
state-and-world invariant, and the per-index package for the processed prefix; each step peels the
lookup, the reader-changing block and the definition builder with the bind inversion and feeds them
to the step hypothesis in that order. The index bookkeeping is the standard append-split reasoning.
\end{proof}

Chaining the invariant through state and world \emph{together} is what lets a per-sibling invariant
be re-established at one sibling from the one held at its predecessor; the cold-start induction of
Chapter~\ref{chap:coldstart} consumes it. Two companions in the same file --- a monotonicity lemma
for the sibling loop, and the transport saying that below the fixpoint the sibling loop's successful
run is verbatim the shipping loop's --- have no consumer in the current tree.

\section{Binder helpers}
\label{sec:run:binders}

\begin{lemma}[The binder helpers]
\label{lem:Erasure-run_withLocalDecl_ok}
\lean{Erasure.run_withLocalDecl_ok, Erasure.run_withLocalDef_ok, Erasure.run_lambdaMonocular_ok,
Erasure.run_letMonocular_ok, Erasure.run_forallMonocular_ok,
Erasure.run_lambdaMonocularOrIntro_ok, Erasure.run_lambdaOrIntroToArity_ok}
\leanok
\uses{lem:Erasure-run_panic, def:Erasure-mkLambda}
A successful run of a continuation-passing binder helper is a successful run of its continuation at
the \emph{same} erasure state, at a fresh identifier, under some extended reader and at some
intermediate world. For the destructuring helpers there is a second possibility: the argument had
the wrong shape, the helper's \code{unreachable!} fired, and --- a panic succeeding --- the run
returned the default value with state and world unchanged. The arity-peeling helper peels exactly as
many binders as the arity asks for.
\end{lemma}

\begin{proof}
\leanok
\uses{lem:Erasure-run_bind_ok, lem:Erasure-run_pure, lem:Erasure-run_panic}
Each helper is a shape test followed either by a fresh-identifier mint, a reader extension and the
continuation, or by \code{unreachable!}. The mint is state-transparent, the reader extension is
\code{withReader}, and the panic arm is closed by Lemma~\ref{lem:Erasure-run_panic}. The
arity-peeling helper is a structural recursion on the arity over the monocular case.
\end{proof}

\begin{remark}
None of these rules constrains the extended reader; that is left to the world-indexed twins of
Chapter~\ref{chap:coldstart}, which additionally report that the continuation's reader has the
\emph{same configuration}. One asymmetry is worth recording: the arity-peeling rule's fall-through
disjunct drops the ``world unchanged'' clause that the other destructuring rules carry, because the
recursive case loses it through an existential. It is forced by the proof, but it weakens the rule
for a world-indexed consumer.
\end{remark}

\begin{lemma}[The \lbox-side binder constructors]
\label{lem:Erasure-run_mkAlt_ok}
\lean{Erasure.run_fvar_to_name_ok, Erasure.run_mkLambda_ok, Erasure.run_mkLetIn_ok,
Erasure.run_mkAlt_ok}
\leanok
\uses{def:toBvar, def:Erasure-mkLambda, def:LBTerm}
Naming a free variable touches neither state nor world: it only reads the reader's local context and
tests the user-facing name for printable ASCII. The \lbox-side binder constructors are otherwise
pure abstraction: a successful \code{mkLambda} run returns a lambda whose body is the input
abstracted over the variable at level zero, and likewise for \code{mkLetIn}. The one with content is
\code{mkAlt}: a case alternative's binder-name list has exactly as many entries as the variables
closed over, and its body is the fold of abstractions over the reversed, index-paired variable list.
\end{lemma}

\begin{proof}
\leanok
\uses{lem:Erasure-run_list_mapM_ok, lem:Erasure-run_list_forIn_ok-prime, lem:Erasure-run_bind_ok}
The variable-naming helper is a read followed by a \code{pure} on either branch of the ASCII test.
\code{mkLambda} and \code{mkLetIn} bind it and then return. \code{mkAlt} runs a \code{mapM} over the
variables for the names and a \code{for} loop for the abstraction fold, both closed by the
corresponding loop rules.
\end{proof}

These realise the locally-nameless-to-de-Bruijn abstraction step that MetaCoq gets for free from its
de Bruijn syntax --- a genuine addition of this development with no counterpart in [S]. The fold
shape is precisely what the \lbox-side closedness lemmas of Chapter~\ref{chap:lowering} compute a
level for.

\section{How the run model is consumed}
\label{sec:run:world}

Everything above is a library; two clients spend it. The refinement of
Chapter~\ref{chap:refinement} runs the family's fixpoint induction with eighteen motives in run-ok
form, each concluding an \Erases-flavoured relation together with the run conclusion
\code{RunConcl}; that is Theorem~\ref{thm:visitExpr_refines_erasesLB}. The cold start of
Chapter~\ref{chap:coldstart} runs it a second time with a \emph{world-indexed} predicate: the
closure bundle \code{RunClosedW} (Definition~\ref{def:RunClosedW}) collects one closure fact for
every way the eighteen-member family can move state or world, and the induction
\code{visitExpr\_shapeW} (Theorem~\ref{thm:visitExpr_shapeW}) pushes such a predicate through the
whole family while concluding, for every member that returns a \lbox{} term, the per-call shape
\code{ShapeCW} (Definition~\ref{def:ShapeCW}): the term is fix-free, de Bruijn closed and in applied
form. Ten of the bundle's clauses are one per ambient primitive the family calls, and they exist
precisely because a state-transparent primitive still advances the name generator, which lives in
the world --- which is why the world-indexed exit rules
(Proposition~\ref{prop:Erasure-run_nonrec_exit_ok-prime}) exist alongside the state-only ones.

Two of the bundle's clauses read the reader's configuration. That is where \code{ConfigPinned}
(Definition~\ref{def:ConfigPinned}) enters: the five configuration restrictions the whole development
is stated under --- no simplification-attribute replacement, no axiomatisation of externally
implemented constants, Peano naturals, no constructor-argument-mask pruning, and no automatic
inlining of typeclass dispatch. Each is a scope restriction stated as a hypothesis rather than an
omission, and the default \code{\#erase} configuration is \textbf{not} the pinned one: it differs on
three of the five fields.

\section{Deviations from the paper and caveats}
\label{sec:run:caveats}

\textbf{No environment pass.} [S, Sec.~7.3] erases a Rocq global environment to a \lbox{}
environment with a separate judgement, before erasing terms against it. This eraser builds the
\lbox{} declarations \emph{during} the term walk, demand-driven and memoised
(Lemma~\ref{lem:Erasure-run_get_constant_kername_ok}). The operational counterpart of
\ErasesEnv's clauses is the registration path of Section~\ref{sec:run:reg}: a constant becomes
either an axiom with no body (Theorem~\ref{thm:Erasure-run_addAxiom_ok}), or a constant declaration
with an erased body, or a member of an emitted inductive block
(Theorem~\ref{thm:Erasure-run_register_inductive_cold_ok}).

\textbf{Program logic, not equational reasoning.} MetaCoq's proofs are inductions over a Rocq
function's own equations; this chapter's vocabulary is Hoare rules for \code{do}-blocks and loops, an
admissibility tower and an approximation transport. That is the price of verifying an eraser that was
not written to be verified.

\textbf{Panics as successful branches.} [S] has no analogue, since MetaCoq's eraser is total. Here
the model must carry a default fall-through everywhere (Lemma~\ref{lem:Erasure-run_panic}), and the
shipping consequences --- the ineffective duplicate guard in \code{addAxiom}, the silently defaulted
binder helpers --- are recorded rather than excluded.

\textbf{No constructor-argument pruning.}
Theorem~\ref{thm:Erasure-run_register_inductive_cold_entries} is where the pipeline's declining of
ConCert-style constructor-argument pruning becomes visible: with the pruning flag off --- one of the
five things \code{ConfigPinned} pins --- every constructor mask is an all-keep mask of exactly the
constructor's field count.

\textbf{Applied-form constructors.} The applied-form convention of this frontend --- a bare
constructor applied through application nodes, rather than the block constructor node MetaCoq and
CertiCoq prefer --- appears here as one of the three output conditions that the world-indexed
induction of Chapter~\ref{chap:coldstart} \emph{concludes} rather than assumes.

\textbf{Two caveats on the sources.} First, the declarations of this chapter live in the root
namespace \code{Erasure}, not in \code{LeanToLambdaBox.Erasure}; the file's module path and its
namespace do not agree. Second, a number of docstrings in \code{ErasureRun.lean} cross-reference
names that no longer exist in the tree, among them the trust class the preprocessing hypothesis is
attributed to and the fixpoint rule of \Erases{} that the principal-argument corollary is said to
serve. No statement is affected, but a reader following such a cross-reference will be misled; the
blueprint deliberately repeats none of those names.
